% =============================================================================
% 中文论文主稿件模板 —— 《计算机辅助设计与图形学学报》投稿格式
%
% 编译：LuaLaTeX（paper/zh/scripts/compile.bat；脚本自动设 OSFONTDIR 指向项目字体）
% 字体：项目内 paper/zh/fonts/（换电脑无需重装系统字体）
%
% 对照期刊格式要点：
%   中文正文 方正书宋_GBK 10pt ｜ 英文/公式/符号 Times New Roman
%   标题 黑体三号 ｜ 作者 方正仿宋四号 ｜ 单位 方正书宋8pt
%   摘要/关键词 方正书宋8.5pt（"摘要/关键词"标签黑体）
%   正文双栏（栏宽21.59字符，栏距3字符）
%   1级标题 11.5pt 黑体 ｜ 2级 五号 黑体 ｜ 3级 书宋10pt
%   参考文献：标题黑体11.5pt，条目8pt（GB/T 7714）
%   注：占位内容用 {花括号} 或直接文字，勿用 [方括号]（会被当作可选参数）
% =============================================================================
\documentclass[UTF8]{ctexart}

% ---------------------------------------------------------------------------
% 宏包
% ---------------------------------------------------------------------------
\usepackage{fontspec}
\usepackage{unicode-math}   % \setmathfont 依赖（含 amssymb 符号，勿再加载 amssymb）
\usepackage{amsmath}
\usepackage{booktabs}       % 三线表
\usepackage{graphicx}
\usepackage{array}
\usepackage{cite}
\usepackage[hidelinks]{hyperref}
\usepackage{xurl}
\usepackage{enumitem}

% ---------------------------------------------------------------------------
% 字体（项目内 paper/zh/fonts/；compile.ps1 设 OSFONTDIR；换电脑无需重装）
% 方正仿宋_GBK 解压到 fonts/FangZhengFangSong-GBK/ 后，把 zhfang 换为对应族名
% ---------------------------------------------------------------------------
\setCJKmainfont{FZShuSong-Z01}      % 方正书宋_GBK（中文正文）
\setCJKsansfont{SimHei}             % 黑体（标题/摘要标签）
\setCJKfamilyfont{zhfang}{FangSong} % 仿宋（作者名；方正仿宋_GBK 解压后换族名）
\setmainfont{Times New Roman}       % 英文/数字/符号
\setmathfont{TeX Gyre Termes Math}  % 公式 Times 风格（TeX Live 自带）

% 三线表线宽（上下 1.0pt，内线 0.5pt）
\setlength{\heavyrulewidth}{1pt}
\setlength{\lightrulewidth}{0.5pt}

\graphicspath{{../figures/final/}}
\setlength{\emergencystretch}{2em}
\Urlmuskip=0mu plus 1mu\relax

% ---------------------------------------------------------------------------
% 中文标题区（单栏通栏）
% ---------------------------------------------------------------------------
\begin{document}

\begin{center}
  % 标题：黑体三号
  {\zihao{3}\heiti {中文标题}}\par
  \vspace{0.6em}
  % 作者：方正仿宋四号（通信作者加 *；同单位不加序号）
  {\zihao{4}\fangsong {作者1}$^{1,2)}$，{作者2}$^{1)*}$}\par
  \vspace{0.3em}
  % 单位：方正书宋 8pt（单位全称到院系级；涉密除外）
  {\fontsize{8pt}{10pt}\selectfont
    1)（{中文单位全称}　{城市}　{邮编}）\quad
    2)（{中文单位全称}　{城市}　{邮编}）}\par
  \vspace{0.3em}
  % 通信作者邮箱
  {\fontsize{8pt}{10pt}\selectfont (E-mail: {邮箱})}\par
\end{center}

\vspace{0.3em}
\noindent{\heiti\fontsize{8.5pt}{11pt}\selectfont 摘\quad 要：}%
{\fontsize{8.5pt}{11pt}\selectfont
{中文摘要：目的(1~2句)、方法过程、实验/算例(数据、对比、评价指标、定量结果)；不用"本文/我们"，用"所提方法/算法/模型"。}}\par

\vspace{0.3em}
\noindent{\heiti\fontsize{8.5pt}{11pt}\selectfont 关键词：}%
{\fontsize{8.5pt}{11pt}\selectfont
{关键词1}；{关键词2}；{关键词3}；{关键词4}；{关键词5}}\par

\vspace{0.3em}
\noindent{\fontsize{8.5pt}{11pt}\selectfont
中图法分类号：{分类号}\hfill DOI：10.3724/SP.J.1089.{稿件编号}}\par

% ---------------------------------------------------------------------------
% 英文标题区（Times New Roman；题目实词首字母大写）
% ---------------------------------------------------------------------------
\begin{center}
  {\fontsize{13.5pt}{16pt}\selectfont\bfseries {English Title}}\par
  \vspace{0.5em}
  {\zihao{5}\selectfont {Author 1}$^{1,2)}$，{Author 2}$^{1)*}$}\par
  \vspace{0.3em}
  {\fontsize{8.5pt}{10pt}\selectfont\itshape
    1) ({English Affiliation}，{City} {Postcode})}\par
  \vspace{0.3em}
  \noindent{\fontsize{9pt}{12pt}\selectfont
  \textbf{Abstract:} {English abstract.}\par
  \textbf{Key words:} {kw1}; {kw2}; {kw3}; {kw4}; {kw5}}
\end{center}

% ---------------------------------------------------------------------------
% 正文：双栏开始
% ---------------------------------------------------------------------------
\twocolumn

\fontsize{10pt}{14pt}\selectfont % 正文 方正书宋 10pt

% ---------- 引言（不加章节号与标题） ----------
{引言：研究背景、动机、与已有工作的区别。不插图、表、公式，不用结论性表述；与摘要、结论不重复。}

% ---------- 1 级标题：11.5pt 黑体 ----------
\section{\heiti\fontsize{11.5pt}{14pt}\selectfont {1级标题}}

{正文。}

% ---------- 2 级标题：五号 黑体 ----------
\subsection{\heiti\zihao{5}\selectfont {2级标题}}

{正文。}

% ---------- 3 级标题：方正书宋 10pt ----------
\subsubsection{\fontsize{10pt}{12pt}\selectfont {3级标题}}

{正文。}

% ---------- 图（示例） ----------
% 图题/分图题精炼、图内不生成标题；图题无句末标点；图可通栏；文字环绕用嵌入型
% 把图文件放 paper/zh/figures/final/ 后取消下面注释
% \begin{figure}[t]
% \centering
% \includegraphics[width=\columnwidth]{fig1_architecture}
% \caption{{图题（图题无标点）}} % 方正书宋 9pt
% \label{fig:architecture}
% \end{figure}

% ---------- 三线表（示例） ----------
% 表题黑体9pt；表序号 Times 加粗；栏头/栏目书宋8pt；三线表(上下1.0pt、内0.5pt)
\begin{table}[t]
\centering
\caption{\heiti\fontsize{9pt}{11pt}\selectfont {表题}} % 黑体9pt
\label{tab:params}
\fontsize{8pt}{10pt}\selectfont
\begin{tabular}{lcc}
\toprule
{栏头} & {横向栏目} & {横向栏目} \\
\midrule
{纵栏目} & 数值 & 数值 \\
{纵栏目} & 数值 & 数值 \\
\bottomrule
\end{tabular}\\[2pt]
\fontsize{8pt}{10pt}\selectfont 注：粗体表示最优值。 % 表注8pt
\end{table}

% ---------- 公式（对应 MathType；被引用的加序号，先给出后引用） ----------
\begin{equation}
  E = m c^2 \label{eq:energy}
\end{equation}
如式~(\ref{eq:energy}) 所示。

% ---------- 算法（输入/输出书宋10pt；Step 书宋9pt；下级首行右移2格） ----------
\textbf{本文算法步骤如下：}\\
输入：{输入描述}。\\
输出：{输出描述}。\\
Step1. {第一步操作}。\\
Step2. {第二步操作}；\\
\hspace{2em} Step2.1. {子步骤，右移2格}；\\
\hspace{2em} Step2.2. {子步骤}。

% ---------- 结语（两字标题加自然空；不重复摘要与引言） ----------
\section{\heiti\fontsize{11.5pt}{14pt}\selectfont 结\quad 语}
{结语：要点概述、优缺点分析、未来展望。}

% ---------- 致谢（非必需） ----------
\section*{\heiti\fontsize{11.5pt}{14pt}\selectfont 致\quad 谢}
{感谢{XX}。非必需项，对不符合署名条件但有贡献者致谢。}

% ---------------------------------------------------------------------------
% 参考文献（标题黑体11.5pt；条目8pt，GB/T 7714，中文文献给英文信息）
% ---------------------------------------------------------------------------
\begin{thebibliography}{99}
\fontsize{8pt}{10pt}\selectfont
\bibitem{ref1} {作者}. {题名}[J]. {期刊名}, {年}, {卷}({期}): {页码}. % 期刊
\bibitem{ref2} {作者}. {题名}[C] //{会议论文集名称}. {出版地}: {出版者}, {年}: {页码}. % 会议
\bibitem{ref3} {作者}. {书名}[M]. {版本}. {出版地}: {出版者}, {年}: {页码}. % 专著
\bibitem{ref4} {作者}. {题名}[D]. {保存地}: {保存单位}, {年}. % 学位论文
\bibitem{ref5} {作者}. {题名}[P]. {国别}, {专利号}. {日期}. % 专利
\bibitem{ref6} {起草者}. {标准代号} {标准名称}[S]. {出版地}: {出版者}, {年}. % 标准
\end{thebibliography}

% ---------------------------------------------------------------------------
% 附录（可选；附录文字9pt）
% ---------------------------------------------------------------------------
% \appendix
% \section*{\heiti\zihao{5}\selectfont 附录A　{标题}}
% \fontsize{9pt}{12pt}\selectfont {附录正文。}

\end{document}
